\section*{Capítulo \ref{ch:g11:forces-and-motion} --- Forças e movimento}

\begin{solution}{exo:g11:forces-and-motion:1}
(a) Velocidade constante: \omterm{def:g11:forces-and-motion:netforce}{força resultante} nula. (b) A velocidade escalar
muda: não nula, oposta ao movimento. (c) A direção muda (a velocidade é
um vetor): não nula, apontando para o interior da curva. (d) Velocidade
constante: nula --- \omterm{def:g10:universal-gravitation:weight}{peso} e arrasto \omterm{def:g10:inertia:compensate}{se compensam}.
\end{solution}

\begin{solution}{exo:g11:forces-and-motion:2}
$d_1/d_2 = 1.0/3.0$ com $d_1 + d_2 = \qty{0.80}{m}$:
$d_1 = \qty{0.20}{m}$ a partir da bola de \qty{3.0}{kg}. Com massas
iguais: no meio, a \qty{0.40}{m} de cada uma.
\end{solution}

\begin{solution}{exo:g11:forces-and-motion:3}
Mesmo $F\Delta t$, massa tripla: o $\Delta v$ do carrinho carregado é
três vezes menor. Para igualar o $\Delta v$ do vazio, empurre com três
vezes mais \omterm{def:g10:inertia:force}{força} (ou por três vezes mais tempo).
\end{solution}

\begin{solution}{exo:g11:forces-and-motion:4}
$\Delta \vect v$: \qty{2}{m/s}, dirigido \emph{para trás} (de \num{10}
para \qty{8}{m/s}). A \omterm{def:g11:forces-and-motion:netforce}{força resultante} aponta no mesmo sentido: oposta a
$\vect v$ --- caso 2, freando.
\end{solution}

\begin{solution}{exo:g11:forces-and-motion:5}
(a) Estático, oposto ao seu empurrão (ainda não há deslizamento).
(b) Cinético, oposto ao deslizamento. (c) Estático, apontando
\emph{ladeira acima} --- contra o deslizamento que a gravidade tenta
iniciar.
\end{solution}

\begin{solution}{exo:g11:forces-and-motion:6}
$\Delta v \propto F\Delta t/m$: (a) \qty{8.0}{m/s}; (b) \qty{2.0}{m/s};
(c) \qty{2.0}{m/s}; (d)
$4.0 \times \frac{3 \times 2}{6} = \qty{4.0}{m/s}$ --- as três mudanças
se cancelam.
\end{solution}

\begin{solution}{exo:g11:forces-and-motion:7}
$v' = \sqrt{8.0^2 + 6.0^2} = \qty{10.0}{m/s}$, com
$\tan\theta = 6.0/8.0$, ou seja, $\theta \approx 37^\circ$ do leste em
direção ao norte. As velocidades se somam como \emph{vetores}:
contribuições perpendiculares se combinam por Pitágoras, e não por soma
de módulos.
\end{solution}

\begin{solution}{exo:g11:forces-and-motion:8}
Já fixados de antemão: a massa da bola e o seu $\Delta v$ (da velocidade
de chegada até zero), e portanto o produto $F\Delta t$. Esticar
$\Delta t$ de \qty{0.02}{s} para \qty{0.08}{s} divide a \omterm{def:g10:inertia:force}{força} média por
$4$.
\end{solution}

\begin{solution}{exo:g11:forces-and-motion:9}
O \omterm{def:g11:forces-and-motion:friction}{atrito estático} da pista sobre a sapatilha: o pé empurra o chão para
trás, e a \omterm{def:g10:inertia:contact}{força de contato} mútua empurra o velocista para a frente. Os
pregos elevam o \omterm{def:g11:forces-and-motion:friction}{atrito estático} máximo (sem escorregar em esforço
total). Sobre gelo sem \omterm{def:g10:inertia:inventory}{atrito} não existe \omterm{def:g10:inertia:force}{força} horizontal externa alguma,
de modo que a velocidade do \omterm{def:g11:forces-and-motion:com}{centro de massa} não pode mudar --- nem
largada, nem caminhada.
\end{solution}

\begin{solution}{exo:g11:forces-and-motion:10}
O ponto de contato de uma roda que gira não desliza: \omterm{def:g11:forces-and-motion:friction}{atrito estático},
mais forte do que o \omterm{def:g11:forces-and-motion:friction}{atrito cinético} de uma derrapagem
(\cref{def:g11:forces-and-motion:friction}) --- frenagem mais curta. Já o
\omterm{def:g11:forces-and-motion:friction}{atrito cinético} de um pneu que derrapa aponta no sentido oposto ao do
deslizamento, seja qual for o ângulo das rodas: girar o volante não muda
nada, e a direção se perde.
\end{solution}

\begin{solution}{exo:g11:forces-and-motion:11}
A tensão do cabo, apontando para dentro, em direção ao atleta. Depois da
soltura, a bola sai em linha reta ao longo da tangente, com a velocidade
que tinha (inércia). O ``cabo'' da Lua é o puxão gravitacional da Terra;
desligue a gravidade e ela parte pela tangente, reta e uniforme, para
sempre.
\end{solution}

\begin{solution}{exo:g11:forces-and-motion:12}
A partir do topo:
$x = \dfrac{1.0 \times 0.60 + 2.0 \times 1.20}{3.0} = \qty{1.00}{m}$.
Distâncias até os centros das duas partes: \qty{0.40}{m} (cabo) e
\qty{0.20}{m} (cabeça), na razão $2 : 1$ --- o inverso da razão de massas
$1.0 : 2.0$, como a \cref{def:g11:forces-and-motion:com} exige.
\end{solution}

\begin{solution}{exo:g11:forces-and-motion:13}
\emph{1.} Mesmos $F$ e $\Delta v$, massa dobrada: a van precisa do dobro
do tempo. \emph{2.} Mesma \omterm{def:g10:relative-motion:speed}{velocidade média} por um tempo dobrado: o dobro
da distância. \emph{3.} $\Delta v$ dobrado com os mesmos $F$ e $m$: o
dobro do tempo; mas a \omterm{def:g10:relative-motion:speed}{velocidade média} também dobra, de modo que a
distância fica multiplicada por $4$. Dobrar a velocidade quadruplica a
distância de frenagem.
\end{solution}

\begin{solution}{exo:g11:forces-and-motion:14}
\emph{1.} $T = 27.3 \times \qty{86400}{s} \approx \qty{2.36e6}{s}$, logo
$v = \dfrac{2\pi \times \num{3.84e8}}{\num{2.36e6}} \approx
\qty{1.0e3}{m/s} \approx \qty{1}{km/s}$. \emph{2.} Em direção à Terra
(para dentro); o puxão gravitacional da Terra. \emph{3.} A Lua
\emph{está} caindo: a variação da velocidade dela aponta para a Terra a
cada instante. A velocidade tangencial apenas garante que ela continue
errando o alvo --- ela cai \emph{em volta} da Terra, em vez de cair
dentro dela.
\end{solution}

\begin{solution}{exo:g11:forces-and-motion:15}
Mesmos $m$ e $\Delta v$, de modo que $F\Delta t$ está fixado: um
$\Delta t$ cinquenta vezes maior na espuma significa uma \omterm{def:g10:inertia:force}{força} média
cinquenta vezes menor. Dobrar os joelhos estica o $\Delta t$ da
aterrissagem de um solavanco para uma flexão longa; um colchão de queda
é, mecanicamente falando, feito de milissegundos extras.
\end{solution}

\begin{solution}{pb:g11:forces-and-motion:1}
\textbf{1.} $50/3.6 \approx \qty{13.9}{m/s}$ (digamos \qty{14}{m/s}).
\textbf{2.} $\Delta v \approx \qty{14}{m/s}$: o ocupante se move a
\qty{14}{m/s} antes e a zero depois --- as duas pontas são fixadas pela
batida, de modo que dispositivo algum toca em $\Delta v$.
\textbf{3.} $\Delta v \propto F\Delta t/m$
(\cref{thm:g11:forces-and-motion:secondlaw}) com $\Delta v$ e $m$ fixos
obriga $F\Delta t$ a ficar fixo. O único botão do projetista é
$\Delta t$.
\textbf{4.} $0.150/0.050 = 3$: a zona de deformação divide a \omterm{def:g10:inertia:force}{força} média
por $3$.
\textbf{5.} Para trás, no sentido oposto ao do movimento --- a velocidade
encolhe de \qty{14}{m/s} até zero.
\textbf{6.} A frente que se amassa é uma máquina para fazer a parada
durar mais: metal que se deforma compra milissegundos.
\textbf{7.} Mesmos $\Delta v$ e $m$, com $\Delta t$ três vezes menor: o
boneco do carro antigo recebe três vezes a \omterm{def:g10:inertia:force}{força}.
\textbf{8.} \omterm{def:g10:relative-motion:speed}{Velocidade média} $\approx \qty{6.9}{m/s}$; distância
$\approx 6.9 \times 0.150 \approx \qty{1.0}{m}$ --- mais ou menos o \omterm{def:g10:orders-of-magnitude:unit}{metro}
de carro deformável que uma zona de deformação real oferece.
\textbf{9.} Um bico de \qty{10}{m} é impossível de dirigir e de
estacionar, e a massa dele pioraria toda batida. A célula dos passageiros
não pode se deformar: é ela que preserva o espaço de sobrevivência.
\textbf{10.} ``\dots então tem de escolher $\Delta t$.''
\textbf{11.} $0.150/0.005 = 30$: a parada no para-brisa é trinta vezes
mais violenta do que a com cinto.
\textbf{12.} Um cinto que estica um pouco alonga mais ainda o $\Delta t$
do próprio ocupante (e espalha a \omterm{def:g10:inertia:force}{força}); um arreio de aço imporia ao
passageiro, mais frágil, a parada mais curta da carroceria.
\textbf{13.} $0.080/0.010 = 8$: o airbag divide a \omterm{def:g10:inertia:force}{força} sobre a cabeça
por $8$.
\textbf{14.} Mesmos $\Delta v$ e $\Delta t$: a \omterm{def:g10:inertia:force}{força} escala com a massa,
$70/20 = 3.5$ vezes maior para o adulto.
\textbf{15.} Variação exigida: $14/0.150 \approx \qty{93}{m/s}$ a cada
segundo --- cerca de $9.5$ vezes o que a gravidade produz
(\qty{9.8}{m/s} por segundo). Os braços precisam puxar cerca de $9.5$
vezes o \omterm{def:g10:universal-gravitation:weight}{peso} do bebê: como segurar uma carga de \qty{95}{kg}. Braço
algum consegue; o bebê é arremessado para a frente.
\textbf{16.} O cinto de um adulto encaminha a \omterm{def:g10:inertia:force}{força} para o quadril e o
tórax; numa criança ele carregaria a barriga e o pescoço. A cadeirinha
aplica a \omterm{def:g10:inertia:force}{força} (menor) a partes fortes do corpo, e o acolchoamento dela
estica ainda mais o $\Delta t$.
\textbf{17.} Se a célula se deforma, os ocupantes perdem o espaço de
sobrevivência e são atingidos pela própria estrutura. Nada deve bater
diretamente na célula rígida: as paradas dos ocupantes ficam a cargo do
cinto e do airbag ($\Delta t$ longo), enquanto a célula ancora as
extremidades que se amassam.
\textbf{18.} Por simetria, os dois carros param no plano de contato: cada
motorista sofre $\Delta v \approx \qty{14}{m/s}$ ao longo de uma parada
com zona de deformação --- essencialmente o teste da parede, e não ``o
dobro da batida''.
\textbf{19.} As \omterm{def:g10:inertia:force}{forças} são iguais e opostas, mas
$\Delta v \propto F\Delta t/m$: o carro pequeno, de massa pequena, sofre
a grande variação de velocidade.
\textbf{20.} Zona de deformação, cinto que estica, airbag e acolchoamento
da cadeirinha --- todos eles esticam a mesma grandeza: $\Delta t$.
\textbf{21.} A batida fixa a massa e a variação de velocidade e,
portanto, o produto \omterm{def:g10:inertia:force}{força} $\times$ tempo: uma colisão suave é uma colisão
longa --- estique o tempo e a \omterm{def:g10:inertia:force}{força} cai na mesma razão.
\end{solution}
